\documentclass[../main.tex]{subfiles}

\begin{document}

\section{Linguagem escolhida e justificativas}

As implementacoes dos metodos numericos foram feitas em Python. A escolha da
linguagem foi motivada pela legibilidade da sintaxe, pela facilidade de separar
o codigo em funcoes pequenas e pela disponibilidade de recursos nativos para
manipular arquivos de texto e arquivos CSV em UTF-8.

Um cuidado importante do projeto foi nao usar bibliotecas prontas para resolver
os metodos numericos. Assim, nao foram usadas funcoes como
\texttt{numpy.polyfit}, rotinas de integracao do \texttt{scipy}, modelos de
regressao do \texttt{sklearn} ou qualquer ferramenta equivalente. As formulas de
minimos quadrados, interpolacao, derivacao, integracao, extrapolacao de
Richardson e quadratura de Gauss foram implementadas diretamente no codigo.

As bibliotecas usadas tiveram apenas papel de apoio:

\begin{itemize}
    \item \texttt{math}, para funcoes elementares como exponencial, logaritmo,
    seno, cosseno e raiz quadrada;
    \item \texttt{csv}, para gravar tabelas de saida em formato CSV;
    \item \texttt{pathlib}, para organizar caminhos de arquivos de forma mais
    clara;
    \item \texttt{matplotlib}, apenas para gerar graficos dos resultados ja
    calculados pelos metodos implementados manualmente.
\end{itemize}

As entradas dos exercicios foram guardadas em arquivos TXT, dentro da pasta de
entradas do projeto. As saidas numericas foram salvas em CSV na pasta de
resultados, e as figuras usadas no relatorio foram salvas em uma pasta propria
para graficos. Essa separacao facilita a repeticao dos testes: a execucao do
script principal de testes regenera os resultados e os graficos do relatorio a
partir dos dados de entrada.

\end{document}
